\documentclass[a4paper, 11pt]{article}
\usepackage[utf8]{inputenc}
\usepackage[T1]{fontenc}
\usepackage[italian]{babel}
\usepackage{geometry}
\usepackage{graphicx}
\usepackage{tabularx}
\usepackage{fancyhdr}
\usepackage{hyperref}

% Margini
\geometry{top=3cm, bottom=3cm, left=2.5cm, right=2.5cm}

% Intestazione e piè di pagina
\pagestyle{fancy}
\fancyhf{}
\fancyhead[L]{\small \textit{0x06tudio} -- Verbale di Riunione}
\fancyhead[R]{\small Anno Accademico 2026/2027}
\fancyfoot[C]{\thepage}
\renewcommand{\headrulewidth}{0.4pt}

\begin{document}
\thispagestyle{plain}

\newcommand{\issue}[1]{\href{https://github.com/0x06tudio/docs/issues/#1}{\textit{#1}}}

\begin{minipage}{0.6\textwidth}
    \vspace{-2cm}
    {\huge \textbf{0x06tudio}} \\[0.5em]
    {\large Progetto didattico di Ingegneria del Software} \\[0.3em]
    {\normalsize Università degli Studi di Padova}
\end{minipage}
\begin{minipage}{0.4\textwidth}
    \raggedleft
    \vspace{0pt}
    \includegraphics[width=6cm]{../../media/logo.png}
    \vspace{1cm}
    \small
    \href{mailto:0x06tudio@gmail.com}{0x06tudio@gmail.com}
\end{minipage}

\vspace{0.5cm}

\begin{center}
    {\LARGE \textbf{Verbale Riunione}}
\end{center}

\vspace{0.5cm}


\begin{tabularx}{\textwidth}{l X}
    \textbf{Versione:} & 0.2.0 \\
    \textbf{Data:} & 2026/10/16 \\
    \textbf{Uso:} & Interno \\
    \textbf{Partecipanti (Presenti):} & Nicolò Beltrame, Pietro Brocadello, Alberto Cavalletto, Catalin Chirilov, Davide Colabove, Andrea Maggio \\
    \textbf{Destinatari:} & 0x06tudio \\
    \textbf{Modalità / Luogo:} & Online / Discord \\
    \textbf{Ora inizio:} & 15:00 \\
    \textbf{Ora fine:} & 17:00 \\
    \textbf{Durata totale:} & 2 ore \\
\end{tabularx}
\section*{Ordine del giorno}
\begin{enumerate}
    \item Definizione way of working\vspace{-0.2cm}
    \item Identità del gruppo\vspace{-0.2cm}
    \item discussione capitolati\vspace{-0.2cm}
\end{enumerate}
\vspace{-0.5cm}

\section*{Registro delle modifiche}
\begin{tabularx}{\textwidth}{|c|c|X|c|c|}
    \hline
    \textbf{Versione} & \textbf{Data} & \textbf{Descrizione} & \textbf{Verificatore} & \textbf{Autore} \\ \hline
    0.2.0 & 2026/11/03 & Correzioni per nuova candidatura &  & Andrea Maggio \\ \hline
    0.1.0 & 2026/10/20 & Creazione e stesura bozza iniziale & Nicolò Beltrame & Andrea Maggio \\ \hline
\end{tabularx}

\newpage
\tableofcontents
\newpage

\section{Resoconto della discussione}

\subsection{Punto 1: Identità del gruppo}
Il gruppo si chiama "0x06tudio", ed ha questo logo:
\newline
\newline
\includegraphics[width=5cm]{../../media/logo.png}
\subsection{Punto 2: Primi aspetti del Way of Working}
Si è deciso di usare LaTeX per la stesura di documenti ufficiali, GitHub per la collaborazione sulla documentazione e sul progetto, git per il versionamento. Si è deciso di effettuare allineamenti interni a cadenza settimanale.
\subsection{Punto 3: Shortlist capitolati}
In seguito ad un sondaggio sulle preferenze di ciascuno, sono state inviate richieste di incontri per gli esponenti dei capitolati C3, C6 e C9.



\section{Riepilogo delle decisioni e degli impegni (To-Do)}

\begin{tabularx}{\textwidth}{|c|X|c|}
    \hline
    \textbf{ID Decisione} & \textbf{Descrizione} & \textbf{Motivazione} \\ \hline
    D 1.1 & Nome 0x06tudio e logo & Esito votazione sondaggio interno \\ \hline
    D 1.2 & Utilizzo di LaTeX & Familiarità da parte dei membri del gruppo \\ \hline
    D 1.3 & Utilizzo di GitHub &  Industry standard \\ \hline
    D 1.4 & Utilizzo di git &  Familiarità da parte dei membri del gruppo \\ \hline
    D 1.5 & Richiesta colloqui in merito ai capitolati C3, C6 e C9 & Necessità di chiarimenti \\  \hline
    D 1.6 & Fissata prossima riunione interna in data 2026/10/23 & Allineamento settimanale \\ \hline
\end{tabularx}

\vspace{1cm}

\noindent
\begin{tabularx}{\textwidth}{|c|X|X|c|c|}
    \hline
    \textbf{Task ID} & \textbf{Descrizione} & \textbf{Assegnato a} & \textbf{Scadenza} & \textbf{Issue} \\ \hline
    TD 1.1 & Creazione recapito mail del gruppo & Davide Colabove & 2026/10/16 & \issue{16} \\ \hline
    TD 1.2 & Creazione logo del gruppo & Alberto Cavalletto & 2026/10/16 & \issue{17} \\ \hline
    TD 1.3 & Invio mail alle aziende & Tutti i membri (tramite recapito del gruppo) & 2026/10/16 & \issue{18} \\ \hline
    TD 1.4 & Creazione della repository GitHub & Tutti i membri & 2026/10/16 & /\\ \hline
    TD 1.5 & Approfondimenti sulle conoscenze richieste per i capitolati & Tutti i membri & 2026/10/19 & \issue{19} \\ \hline
\end{tabularx}

\end{document}